\section{Analytical Investigation of Crossing Sequence Memory}

\subsection{Preliminaries}

Let
\begin{equation}
z_i=e^{\mathrm{i}\theta_i},
\qquad
\xi_i^\mu=e^{\mathrm{i}\phi_i^\mu},
\qquad
i=1,\ldots,N,
\end{equation}
and define the normalized complex overlap
\begin{equation}
m_\mu(z)
=
\frac{1}{N}
\sum_{j=1}^N
\overline{\xi_j^\mu}z_j.
\end{equation}
The autonomous Kuramoto dynamics are
\begin{equation}
\dot{\theta}_i
=
g,\operatorname{Im}
\left[
h_i(z)\overline{z_i}
\right].
\end{equation}
Equivalently, since $\dot z_i=\mathrm{i}z_i\dot\theta_i$,
\begin{equation}
\dot z_i
=
\frac{g}{2}
\left[
h_i(z)
-
z_i^2\overline{h_i(z)}
\right].
\label{eq:z-dynamics}
\end{equation}

We consider fields that may be written in the form
\begin{equation}
h_i(z)
=
\sum_\nu
\xi_i^\nu A_\nu(m,\overline m),
\label{eq:general-memory-field}
\end{equation}
where each $A_\nu$ is a static polynomial of bidegree $(2,1)$ in the overlaps.

For example, a transition $\mu\rightarrow\mu+1$ has coefficients
\begin{align}
A_{\mu+1}^{\mathrm{M1}}
&=
m_\mu |m_\mu|^2,
\
A_{\mu+1}^{\mathrm{B}}
&=
m_\mu |m_{\mu-1}|^2,
\
A_{\mu+1}^{\mathrm{CB}}
&=
m_\mu^2\overline{m_{\mu-1}},
\
A_{\mu+1}^{\mathrm{CT}}
&=
m_\mu m_{\mu-2}\overline{m_{\mu-1}}.
\end{align}

Throughout the analysis it is useful to distinguish two notions of context.

\begin{definition}[Explicit dynamical context]
A model has explicit dynamical context if its state contains a variable in
addition to $z$, for example a delayed trace $q$, such that two trajectories
with the same instantaneous oscillator state $z$ may nevertheless have
different future evolution because their additional states differ.
\end{definition}

\begin{definition}[Residual geometric context]
A static rule has residual geometric context if branch information is encoded
only through small differences in the instantaneous state $z$, and hence in
$m(z)$, as the trajectory approaches a nominally shared pattern.
\end{definition}

The bigram, coherent-bigram, and coherent-trigram rules considered here belong
to the second category.  They do not violate the autonomous current-state
obstruction derived below.

\subsection{No-go result for exact crossings}

Consider two stored sequences
\begin{equation}
\cdots
\rightarrow \xi^{p_1}
\rightarrow \xi^\star
\rightarrow \xi^{s_1}
\rightarrow\cdots
\end{equation}
and
\begin{equation}
\cdots
\rightarrow \xi^{p_2}
\rightarrow \xi^\star
\rightarrow \xi^{s_2}
\rightarrow\cdots,
\end{equation}
with $s_1\neq s_2$, but with the central state $\xi^\star$ shared exactly.

\begin{proposition}[Exact-crossing obstruction]
\label{prop:no-go}
Let the network be deterministic and autonomous, with
\begin{equation}
h(z)=H(m(z),\overline{m(z)})
\end{equation}
for a single-valued static function $H$.  If two trajectories arrive at exactly
the same oscillator state
\begin{equation}
z=\xi^\star,
\end{equation}
then they have identical instantaneous vector fields at the crossing,
independent of the branch from which they arrived.
Consequently, exact deterministic branch selection cannot depend on the
incoming branch once the complete dynamical state has collapsed to
$\xi^\star$.
\end{proposition}

\paragraph{Proof.}
At $z=\xi^\star$, every overlap is uniquely determined:
\begin{equation}
m_\mu^\star
=
m_\mu(\xi^\star)
=
\frac{1}{N}
\sum_j
\overline{\xi_j^\mu}\xi_j^\star.
\end{equation}
This vector is a function only of $\xi^\star$ and the stored patterns.  It does
not depend on the trajectory by which $\xi^\star$ was reached.  Therefore
\begin{equation}
h(\xi^\star)
=
H(m^\star,\overline{m^\star})
\end{equation}
is identical for all incoming branches.  The deterministic vector field
\begin{equation}
\dot{\theta}_i
=
g\operatorname{Im}
\left[
h_i(\xi^\star)\overline{\xi_i^\star}
\right]
\end{equation}
is consequently identical as well.  Hence branch-dependent continuation from
the exactly identical complete state is impossible. \hfill$\square$

\paragraph{Corollary.}
The obstruction applies not only to Markov-1 rules, but to every static
polynomial rule constructed from the instantaneous overlaps, including the
bigram, coherent-bigram, coherent-trigram, multilag, and ridge-fitted rules.

Thus, when these rules empirically resolve crossings, the relevant mechanism is
not true hidden memory.  Rather, the trajectory fails to collapse exactly onto
$\xi^\star$ before leaving its neighborhood, and a small branch-dependent
component survives in the overlap vector.

Exact deterministic disambiguation at an identical state would require at
least one of the following:
\begin{equation}
\text{additional dynamical state},\qquad
\text{external branch bias},\qquad
\text{explicit delay},\qquad
\text{stochastic symmetry breaking},
\end{equation}
or a representation in which the apparently shared state is lifted into
distinct context-dependent states.

\subsection{Local overlap expansion near a crossing}

We next analyze the empirically relevant case in which the trajectory passes
near, rather than exactly through, the shared pattern.

Because $z_i$ must remain on the unit circle, the schematic expression
\begin{equation}
z
\approx
\xi^\star+\varepsilon\xi^{p}
\end{equation}
should be interpreted as an overlap-coordinate expansion rather than a literal
linear superposition of phase vectors.

Suppose the trajectory approaching $\xi^\star$ from branch $b$ has overlaps
\begin{align}
m_\star
&=
1-O(\varepsilon^2)-O_p(N^{-1/2}),
\label{eq:m-star-local}
\
m_{p_b}
&=
q_\rho+\varepsilon r_b+\eta_b,
\label{eq:m-prev-local}
\
m_{p_c}
&=
q_\rho+\varepsilon r_{bc}+\eta_c,
\qquad c\neq b.
\label{eq:m-wrong-prev-local}
\end{align}
Here
\begin{equation}
q_\rho
=
\frac{1}{N}
\sum_i
\overline{\xi_i^{p_b}}\xi_i^\star
\end{equation}
is the typical predecessor–crossing overlap induced by the corruption
procedure.

For the benchmark distance
\begin{equation}
d(\xi,\zeta)
=
1-
\left|
\frac{1}{N}
\sum_i
\overline{\xi_i}\zeta_i
\right|,
\end{equation}
one expects schematically
\begin{equation}
|q_\rho|
\simeq
1-\rho,
\end{equation}
up to details of the corruption distribution.

The variable $r_b$ measures alignment between the surviving perturbation and
the incoming branch direction.  For the correct branch we assume
\begin{equation}
r_b=O(\sigma_\rho),
\end{equation}
whereas independent branch directions give
\begin{equation}
r_{bc}
=
O_p\left(\frac{\sigma_\rho}{\sqrt N}\right),
\qquad c\neq b.
\end{equation}
A natural scale is
\begin{equation}
\sigma_\rho^2
\sim
1-|q_\rho|^2,
\end{equation}
which measures the innovation orthogonal to the shared pattern.

The finite-size overlap fluctuations satisfy
\begin{equation}
\eta_b,\eta_c
=
O_p(N^{-1/2})
\end{equation}
under a random-phase approximation.

For fully independent random phase patterns, $q_\rho=0$, giving the simpler
idealization
\begin{equation}
m_{p_b}
=
\varepsilon r_b+O_p(N^{-1/2}),
\qquad
m_{p_c}
=
O_p(N^{-1/2}).
\end{equation}

The corruption-path benchmarks are more structured than this independent
pattern limit, because all immediate predecessors of a crossing may have a
nonzero common overlap $q_\rho$ with the crossing pattern.

\subsection{Expansion of the branch-selecting fields}

Suppose there are $B$ candidate transitions
\begin{equation}
p_b\rightarrow \star\rightarrow s_b,
\qquad
b=1,\ldots,B.
\end{equation}

At a state near $\xi^\star$, the part of each rule that points toward successor
$s_b$ can be studied through its coefficient.

\subsubsection{Markov-1 rule}

For all branches,
\begin{equation}
A_{s_b}^{\mathrm{M1}}
=
m_\star|m_\star|^2.
\end{equation}
Hence
\begin{equation}
A_{s_1}^{\mathrm{M1}}
=
\cdots
=
A_{s_B}^{\mathrm{M1}}
\end{equation}
whenever the shared pattern is indexed identically as the current state.

There is therefore no branch-selective term even slightly away from the
crossing unless branch information also enters through some additional
feature of the rule.

This is the local manifestation of Proposition~\ref{prop:no-go}.

\subsubsection{Magnitude-gated bigram rule}

For branch $b$,
\begin{equation}
A_{s_b}^{\mathrm{B}}
=
m_\star |m_{p_b}|^2.
\end{equation}
Using $m_\star=1+O(\varepsilon^2)$,
\begin{equation}
A_{s_b}^{\mathrm{B}}
=
\left|
q_\rho+\varepsilon r_b+\eta_b
\right|^2
+
O(\varepsilon^2|m_{p_b}|^2).
\end{equation}
Expanding,
\begin{equation}
A_{s_b}^{\mathrm{B}}
=
|q_\rho|^2
+
2\varepsilon
\operatorname{Re}
\left(
\overline{q_\rho}r_b
\right)
+
\varepsilon^2|r_b|^2
+
O_p(N^{-1/2}).
\label{eq:bigram-expansion}
\end{equation}

The common term $|q_\rho|^2$ does not select a branch.  The leading
branch-specific contrast is therefore
\begin{equation}
S_{\mathrm{B}}
\sim
2\varepsilon
\operatorname{Re}
\left(
\overline{q_\rho}r_b
\right)
\end{equation}
when $q_\rho\neq0$.

In the independent-pattern limit $q_\rho=0$, the linear term vanishes and
\begin{equation}
S_{\mathrm{B}}
\sim
\varepsilon^2|r_b|^2.
\label{eq:bigram-independent-signal}
\end{equation}

Thus magnitude gating may produce a quadratic branch signal in the
uncorrelated limit, but a linear differential signal when the predecessor and
crossing state possess a substantial baseline correlation.

\subsubsection{Coherent bigram rule}

The successor coefficient is
\begin{equation}
A_{s_b}^{\mathrm{CB}}
=
m_\star^2\overline{m_{p_b}}.
\end{equation}
Hence
\begin{equation}
A_{s_b}^{\mathrm{CB}}
=
\overline{q_\rho}
+
\varepsilon\overline{r_b}
+
\overline{\eta_b}
+
O(\varepsilon^2).
\label{eq:coherent-bigram-expansion}
\end{equation}

The branch-independent baseline $\overline{q_\rho}$ again produces no
disambiguating information.  The branch-dependent increment is linear:
\begin{equation}
S_{\mathrm{CB}}
\sim
\varepsilon |r_b|.
\label{eq:coherent-bigram-signal}
\end{equation}

This is one reason the coherent rule may have a stronger local
signal-to-crosstalk ratio than a magnitude-gated rule when the residual context
is small.

\subsubsection{Coherent trigram rule}

For a transition $\mu\rightarrow\mu+1$,
\begin{equation}
A_{\mu+1}^{\mathrm{CT}}
=
m_\mu m_{\mu-2}\overline{m_{\mu-1}}.
\end{equation}

For a one-node crossing, the two lagged overlaps may both retain branch
dependence.  Writing
\begin{equation}
m_{\mu-1}
=
q_1+\varepsilon r_1+\eta_1,
\qquad
m_{\mu-2}
=
q_2+\varepsilon r_2+\eta_2,
\end{equation}
one obtains
\begin{align}
A_{\mu+1}^{\mathrm{CT}}
&=
m_\mu
\left(q_2+\varepsilon r_2+\eta_2\right)
\left(
\overline{q_1}
+
\varepsilon\overline{r_1}
+
\overline{\eta_1}
\right)
\
&=
m_\mu q_2\overline{q_1}
+
\varepsilon m_\mu
\left(
r_2\overline{q_1}
+
q_2\overline{r_1}
\right)
+
O(\varepsilon^2)
+
O_p(N^{-1/2}).
\end{align}

The product therefore contains phase-sensitive information about the relative
geometry of two recent patterns.  This becomes especially useful when one
lagged pattern is shared and the other is still branch-specific.

\subsection{Signal-to-crosstalk scaling}

We now give heuristic mean-field estimates for the branch signal.

These estimates are not rigorous because the overlaps become dynamically
correlated during retrieval.  They are intended as a scaling theory against
which simulations can be compared.

Let
\begin{equation}
P
\end{equation}
denote the number of stored transitions and
\begin{equation}
B
\end{equation}
the number of candidate successors at the crossing.

For unrelated random phase patterns,
\begin{equation}
\frac{1}{N}
\sum_i
\overline{\xi_i^\alpha}\xi_i^\beta
=
O_p(N^{-1/2}),
\qquad
\alpha\neq\beta.
\end{equation}

Accordingly, an inactive overlap has typical magnitude
\begin{equation}
|m_\alpha|
\sim
N^{-1/2}.
\end{equation}

\subsubsection{Bigram scaling}

For the correct branch,
\begin{equation}
S_{\mathrm{B}}
\sim
\begin{cases}
\varepsilon^2\sigma_\rho^2,
& q_\rho\simeq0,
\[4pt]
\varepsilon |q_\rho|\sigma_\rho,
& |q_\rho|>0.
\end{cases}
\end{equation}

For an incorrect branch at the crossing in the independent-pattern limit,
\begin{equation}
|A_{s_c}^{\mathrm{B}}|
\sim
|m_\star||m_{p_c}|^2
\sim
N^{-1}.
\end{equation}
Combining $B-1$ approximately independent competitors gives the scale
\begin{equation}
C_{\mathrm{wrong}}^{\mathrm{B}}
\sim
\frac{\sqrt B}{N}.
\end{equation}

A generic inactive transition has
\begin{equation}
|m_\mu||m_{\mu-1}|^2
\sim
N^{-3/2},
\end{equation}
so $P$ incoherent transitions contribute approximately
\begin{equation}
C_{\mathrm{noise}}^{\mathrm{B}}
\sim
\frac{\sqrt P}{N^{3/2}}.
\end{equation}

A simple sufficient scaling condition is therefore
\begin{equation}
S_{\mathrm{B}}
\gg
c_1\frac{\sqrt B}{N}
+
c_2\frac{\sqrt P}{N^{3/2}},
\label{eq:bigram-condition}
\end{equation}
where $c_1,c_2=O(1)$ depend on the desired error probability and pattern
statistics.

In the uncorrelated case this gives roughly
\begin{equation}
\varepsilon
\gg
\frac{B^{1/4}}
{\sqrt{N},\sigma_\rho}
\end{equation}
when branch competition dominates.

\subsubsection{Coherent-bigram scaling}

The coherent bigram has correct differential signal
\begin{equation}
S_{\mathrm{CB}}
\sim
\varepsilon\sigma_\rho.
\end{equation}

An incorrect candidate branch has
\begin{equation}
|A_{s_c}^{\mathrm{CB}}|
\sim
N^{-1/2},
\end{equation}
and hence
\begin{equation}
C_{\mathrm{wrong}}^{\mathrm{CB}}
\sim
\sqrt{\frac{B}{N}}.
\end{equation}

Inactive bulk terms satisfy
\begin{equation}
|m_\mu|^2|m_{\mu-1}|
\sim
N^{-3/2},
\end{equation}
giving
\begin{equation}
C_{\mathrm{noise}}^{\mathrm{CB}}
\sim
\frac{\sqrt P}{N^{3/2}}.
\end{equation}

Thus a sufficient scaling condition is
\begin{equation}
\varepsilon\sigma_\rho
\gg
c_3\sqrt{\frac{B}{N}}
+
c_4\frac{\sqrt P}{N^{3/2}}.
\label{eq:cb-condition}
\end{equation}

The coherent-bigram signal is linear in $\varepsilon$, but the branch
competitor term is also larger than for magnitude gating.  Which rule performs
better therefore depends on the finite-$N$ geometry and on the common
correlation $q_\rho$.

\subsubsection{Effect of corruption distance}

Introduce
\begin{equation}
c_\rho
:=
|q_\rho|,
\qquad
\sigma_\rho
:=
\sqrt{1-c_\rho^2}.
\end{equation}

For the overlap-distance construction,
\begin{equation}
c_\rho
\approx
1-\rho
\end{equation}
is a natural first approximation.

Very small $\rho$ gives
\begin{equation}
c_\rho\approx1,
\qquad
\sigma_\rho\ll1.
\end{equation}
Successive memories are then very similar, so the innovation distinguishing
branches is weak.

At very large $\rho$, neighboring patterns are easier to separate, but the
trajectory may also lose useful residual predecessor overlap more rapidly.

The simplest qualitative expectation is therefore that branch resolution is
controlled not directly by $\rho$ alone, but by a product such as
\begin{equation}
\varepsilon(\rho)\sigma_\rho,
\end{equation}
where $\varepsilon(\rho)$ measures how much predecessor information survives
when the trajectory reaches the crossing neighborhood.

A generic branch-resolution condition may consequently be written
\begin{equation}
S_{\mathrm{branch}}
>
C_{\mathrm{wrong}}
+
C_{\mathrm{noise}},
\label{eq:generic-snr-condition}
\end{equation}
with
\begin{equation}
S_{\mathrm{branch}}
=
S_{\mathrm{branch}}
\left(
\varepsilon,
\rho
\right)
\end{equation}
and the noise terms scaling with $B$, $P$, and $N$ as above.

\subsection{Shared segments and context lifetime}

Consider now sequences
\begin{align}
p_b
\rightarrow
s_1
\rightarrow
s_2
\rightarrow\cdots\rightarrow
s_L
\rightarrow
y_b,
\qquad b=1,\ldots,B,
\end{align}
where
\begin{equation}
s_1,\ldots,s_L
\end{equation}
are exactly shared.

Let
\begin{equation}
r_k
\end{equation}
denote the branch-specific residual after traversing $k$ shared patterns.

A simple phenomenological approximation is
\begin{equation}
r_{k+1}
\simeq
\lambda_k r_k+\zeta_k,
\end{equation}
where $|\lambda_k|<1$ represents contraction toward the common shared
trajectory and $\zeta_k$ denotes finite-size/crosstalk forcing.

If the shared segment is approximately homogeneous,
\begin{equation}
\lambda_k\simeq\lambda,
\end{equation}
then
\begin{equation}
|r_L|
\simeq
|\lambda|^L |r_0|
\end{equation}
before accounting for accumulated noise.

If the relevant crosstalk threshold is
\begin{equation}
r_{\mathrm{crit}}
\sim
c_B\sqrt{\frac{B}{N}}
+
c_P\frac{\sqrt P}{N^{3/2}},
\end{equation}
a necessary heuristic condition for residual-context disambiguation is
\begin{equation}
|\lambda|^L |r_0|
\gtrsim
r_{\mathrm{crit}}.
\label{eq:context-lifetime}
\end{equation}

Equivalently,
\begin{equation}
L
\lesssim
\frac{
\log(|r_0|/r_{\mathrm{crit}})
}{
-\log|\lambda|
}.
\label{eq:max-shared-length}
\end{equation}

This defines an effective context lifetime.

\subsubsection{Why one-step bigram context fails for longer shared segments}

Consider
\begin{equation}
p_b
\rightarrow
s_1
\rightarrow
s_2
\rightarrow
y_b.
\end{equation}

At the transition out of $s_1$, the bigram factor
\begin{equation}
|m_{p_b}|^2
\end{equation}
is branch-specific.

At the transition out of $s_2$, however, its explicitly indexed lag is
$s_1$, which is identical across branches:
\begin{equation}
A_{y_b}^{\mathrm{B}}
=
m_{s_2}|m_{s_1}|^2.
\end{equation}
The formal one-step context term is therefore no longer branch-specific.

A bigram may still succeed empirically if the current state has not collapsed
onto the shared subspace and $m_{s_1}$ itself carries branch-dependent
contamination.  Such success relies on analog residual dynamics rather than
explicit access to the unique predecessor.

\subsubsection{Why the coherent trigram can resolve a two-pattern shared segment}

For
\begin{equation}
p_b
\rightarrow
s_1
\rightarrow
s_2
\rightarrow
y_b,
\end{equation}
the coherent-trigram coefficient at the exit is
\begin{equation}
A_{y_b}^{\mathrm{CT}}
=
m_{s_2}
m_{p_b}
\overline{m_{s_1}}.
\label{eq:ct-shared-two}
\end{equation}

Here $s_1$ is shared, but $p_b$ is branch-specific.  Hence the factor
\begin{equation}
m_{p_b}
\end{equation}
provides an explicit two-step branch cue.

If
\begin{equation}
m_{s_2}=O(1),
\qquad
m_{s_1}=O(1),
\qquad
m_{p_b}=r_2,
\end{equation}
then
\begin{equation}
S_{\mathrm{CT}}
\sim
|r_2|.
\end{equation}

Thus the coherent trigram naturally matches a shared segment of length two.

For $L>2$, however, the explicitly indexed pattern $\mu-2$ also lies inside
the shared segment.  The coherent trigram therefore does not provide an
arbitrarily long discrete context memory.

For a general shared segment of length $L$, an explicitly indexed context rule
would need access to a lag of at least $L$, or an additional dynamical trace
whose lifetime satisfies Eq.~\eqref{eq:context-lifetime}.

Within a fixed fourth-order $(2,1)$ polynomial budget, only a limited number of
lagged overlap factors can appear simultaneously.  This imposes a structural
constraint on explicit context depth.

\subsection{Reduced overlap dynamics}

The oscillator dynamics admit an exact overlap equation that is useful as a
starting point for mean-field analysis.

Define the pattern Gram matrix
\begin{equation}
G_{\alpha\nu}
=
\frac{1}{N}
\sum_i
\overline{\xi_i^\alpha}\xi_i^\nu.
\end{equation}

Also define the state-dependent second-harmonic tensor
\begin{equation}
R_{\alpha\nu}(z)
=
\frac{1}{N}
\sum_i
\overline{\xi_i^\alpha}
\overline{\xi_i^\nu}
z_i^2.
\label{eq:R-definition}
\end{equation}

Using Eq.~\eqref{eq:z-dynamics},
\begin{align}
\dot m_\alpha
&=
\frac{1}{N}
\sum_i
\overline{\xi_i^\alpha}\dot z_i
\
&=
\frac{g}{2N}
\sum_i
\overline{\xi_i^\alpha}
\left(
h_i-z_i^2\overline{h_i}
\right).
\end{align}

Substituting
\begin{equation}
h_i
=
\sum_\nu
\xi_i^\nu A_\nu,
\end{equation}
gives the exact identity
\begin{equation}
\boxed{
\dot m_\alpha
=
\frac{g}{2}
\left[
\sum_\nu
G_{\alpha\nu}A_\nu
-
\sum_\nu
R_{\alpha\nu}(z)
\overline{A_\nu}
\right].
}
\label{eq:exact-overlap-dynamics}
\end{equation}

Equation~\eqref{eq:exact-overlap-dynamics} is exact, but it is not closed in
the overlaps because $R_{\alpha\nu}(z)$ contains second harmonics of the
microscopic phase state.

\subsection{Random-pattern mean-field approximation}

For approximately orthogonal random phase patterns,
\begin{equation}
G_{\alpha\nu}
=
\delta_{\alpha\nu}
+
O_p(N^{-1/2}).
\end{equation}

If the state is close to a single pattern $\xi^\gamma$,
\begin{equation}
z_i
\approx
\xi_i^\gamma,
\end{equation}
then
\begin{equation}
R_{\alpha\nu}
\approx
\frac{1}{N}
\sum_i
\overline{\xi_i^\alpha}
\overline{\xi_i^\nu}
(\xi_i^\gamma)^2.
\end{equation}

For independent random phases,
\begin{equation}
R_{\gamma\gamma}
\approx1,
\end{equation}
whereas generic combinations
\begin{equation}
(\alpha,\nu)\neq(\gamma,\gamma)
\end{equation}
are of order
\begin{equation}
O_p(N^{-1/2}).
\end{equation}

Thus, for an overlap $\alpha$ transverse to the currently dominant pattern,
one obtains the leading approximation
\begin{equation}
\dot m_\alpha
\approx
\frac{g}{2}A_\alpha
+
O_p
\left(
\frac{|A|}{\sqrt N}
\right).
\label{eq:transverse-overlap}
\end{equation}

This provides a simple interpretation: to leading order, the polynomial
coefficient multiplying pattern $\xi^\alpha$ is also the driving term for the
growth of its overlap.

For the dominant pattern itself, the second term in
Eq.~\eqref{eq:exact-overlap-dynamics} cannot be neglected because
\begin{equation}
R_{\gamma\gamma}=O(1).
\end{equation}
It supplies the nonlinear phase-normalization effect associated with the
constraint $|z_i|=1$.

\subsection{Reduced equations for the crossing modes}

Near a crossing, retain only
\begin{equation}
x=m_\star,
\qquad
r_b=m_{p_b},
\qquad
y_b=m_{s_b}.
\end{equation}

For the Markov rule,
\begin{equation}
A_{s_b}
=
x|x|^2,
\end{equation}
so Eq.~\eqref{eq:transverse-overlap} gives
\begin{equation}
\dot y_b
\approx
\frac{g}{2}
x|x|^2
+
\eta_b^{\mathrm{M1}}.
\end{equation}
All candidate successor modes receive the same leading drive.

For the magnitude-gated bigram,
\begin{equation}
\dot y_b
\approx
\frac{g}{2}
x|r_b|^2
+
\eta_b^{\mathrm{B}}.
\label{eq:reduced-bigram}
\end{equation}

For the coherent bigram,
\begin{equation}
\dot y_b
\approx
\frac{g}{2}
x^2\overline{r_b}
+
\eta_b^{\mathrm{CB}}.
\label{eq:reduced-cb}
\end{equation}

For the coherent trigram, introducing
\begin{equation}
q_b=m_{\mu-2}^{(b)},
\end{equation}
gives
\begin{equation}
\dot y_b
\approx
\frac{g}{2}
x q_b\overline{r_b}
+
\eta_b^{\mathrm{CT}}.
\label{eq:reduced-ct}
\end{equation}

The noise terms summarize Gram-matrix crosstalk, neglected second-harmonic
terms, and all inactive transitions.

Equations~\eqref{eq:reduced-bigram}–\eqref{eq:reduced-ct} provide a minimal
reduced dynamical description of branch competition.

\subsection{Branch-contrast equations}

A particularly useful quantity is the difference between the correct successor
overlap and an incorrect successor overlap:
\begin{equation}
\Delta_{bc}
=
y_b-y_c.
\end{equation}

For the Markov rule,
\begin{equation}
\dot\Delta_{bc}
\approx
\eta_b-\eta_c,
\end{equation}
so there is no deterministic leading-order branch signal.

For the magnitude-gated bigram,
\begin{equation}
\dot\Delta_{bc}
\approx
\frac{g}{2}
x
\left(
|r_b|^2-|r_c|^2
\right)
+
\eta_{bc}.
\label{eq:contrast-bigram}
\end{equation}

For the coherent bigram,
\begin{equation}
\dot\Delta_{bc}
\approx
\frac{g}{2}
x^2
\left(
\overline{r_b}
-
\overline{r_c}
\right)
+
\eta_{bc}.
\label{eq:contrast-cb}
\end{equation}

For the coherent trigram,
\begin{equation}
\dot\Delta_{bc}
\approx
\frac{g}{2}
x
\left(
q_b\overline{r_b}
-
q_c\overline{r_c}
\right)
+
\eta_{bc}.
\label{eq:contrast-ct}
\end{equation}

These equations expose the qualitative distinction between the rules:
Markov-1 contains no branch contrast; the bigram uses a difference of context
magnitudes; the coherent bigram uses the complex predecessor overlap itself;
and the coherent trigram uses a phase-sensitive product of two lagged
overlaps.

\subsection{A heuristic branch-selection proposition}

\begin{proposition}[Heuristic local branch-selection condition]
Suppose a crossing neighborhood is entered from branch $b$ and the reduced
successor dynamics satisfy
\begin{equation}
\dot y_b
=
S_b+\eta_b,
\qquad
\dot y_c
=
S_c+\eta_c,
\end{equation}
for $c\neq b$.  If, throughout the branch-selection interval,
\begin{equation}
\operatorname{Re}
\left[
S_b-S_c
\right]
>
\Gamma_{bc},
\label{eq:branch-sufficient}
\end{equation}
where $\Gamma_{bc}$ bounds deterministic crosstalk and stochastic finite-size
fluctuations in the relevant successor direction, then the correct successor
overlap develops a positive deterministic advantage over competitor $c$.
\end{proposition}

This is not yet a theorem guaranteeing global convergence to $s_b$, because
later nonlinear dynamics may alter the ordering.  It is instead a local
sufficient condition for the correct successor mode to begin winning the
competition.

Using the preceding mean-field estimates, Eq.~\eqref{eq:branch-sufficient}
takes the schematic form
\begin{equation}
S_{\mathrm{branch}}
\gtrsim
C_{\mathrm{wrong}}
+
C_{\mathrm{bulk}}
+
C_{\mathrm{finite}\text{-}N}.
\end{equation}

\subsection{Predictions for numerical experiments}

The analysis suggests several directly testable finite-size predictions.

\begin{enumerate}
\item At an artificially reset exact crossing state
\begin{equation}
z=\xi^\star,
\end{equation}
no deterministic static rule considered here should retain information
about which branch was previously followed.

\item If trajectories are instead allowed to approach the crossing
naturally, branch success should correlate with a measurable residual
contrast such as
\begin{equation}
    |m_{p_b}|-|m_{p_c}|
\end{equation}
or its complex phase-sensitive analogue.
\item Increasing $N$ should reduce random overlap crosstalk approximately
as powers of $N^{-1/2}$, while the deterministic residual signal need not
vanish with $N$.  Consequently, success should sharpen with increasing
system size when the residual context remains $O(1)$.
\item Increasing the number of branches should increase the maximum
competitor fluctuation roughly with a $\sqrt{B}$ scaling at the simplest
rms level, with possible logarithmic corrections if one studies the
maximum over many independent competitors.
\item Increasing the total number of stored transitions should raise bulk
crosstalk approximately as $\sqrt P$ under incoherent-sum assumptions.
\item For shared segments, the branch-specific residual should decay as the
shared segment length increases.  Plotting the measured residual against
$L$ should allow an effective contraction factor $\lambda$ to be estimated.
\item A one-step bigram should deteriorate sharply when the explicit
predecessor entering its gate becomes a shared state, whereas the coherent
trigram should remain effective for a two-pattern shared segment because it
explicitly references the last unique predecessor.
\item For shared segments longer than the rule's explicit lag depth, any
remaining success must arise from analog residual contamination of the
instantaneous state rather than from exact discrete lag indexing.

\end{enumerate}

\subsection{Interpretation}

The analytical picture is therefore a near-crossing, rather than
exact-crossing, mechanism.

At an exact shared state, autonomous deterministic current-state dynamics
cannot know the incoming branch.  This yields a structural no-go result that is
independent of the polynomial order of the static learning rule.

Near the shared state, however, the oscillator trajectory generally retains a
small branch-dependent residual.  In overlap coordinates this appears as
\begin{equation}
m_{p_b}
=
\text{common component}
+
\text{branch-specific residual}
+
\text{finite-size crosstalk}.
\end{equation}
Context-sensitive polynomial rules transform this otherwise small residual
into a differential drive on the candidate successors.

The Markov-1 rule has no such differential term.  The magnitude-gated bigram
uses
\begin{equation}
|m_{p_b}|^2,
\end{equation}
the coherent bigram uses
\begin{equation}
\overline{m_{p_b}},
\end{equation}
and the coherent trigram uses the phase-sensitive two-lag product
\begin{equation}
m_{\mu-2}\overline{m_{\mu-1}}.
\end{equation}

The resulting branch-selection problem is naturally governed by a
signal-to-crosstalk inequality
\begin{equation}
S_{\mathrm{branch}}
>
C_{\mathrm{wrong}}
+
C_{\mathrm{noise}}.
\end{equation}
For random phase patterns, crosstalk decreases with system size, whereas the
useful signal is controlled primarily by the magnitude and lifetime of the
residual context.  For a shared segment of length $L$, reliable disambiguation
requires this residual to remain above the finite-size crosstalk floor until
the branch point is reached.

This framework explains why a coherent trigram can be the simplest
interpretable fourth-order rule that succeeds on a two-pattern shared segment:
it does not evade the exact-state no-go theorem, but it can exploit a
branch-specific overlap two steps in the past while that information is still
present in the instantaneous phase configuration.

